% ================================================================
%  style/preamble.tex
%  《现代催化研究方法精要：原理、读图与方法选择》
%  版式与宏包统一配置 —— 依据企划 v0.1 §6
%  引擎：XeLaTeX      构建：latexmk -xelatex -outdir=build
%
%  【修改纪律】本文件是全书的版式契约。任何改动都会影响 12 章。
%  改动后必须：(1) 重编译样章；(2) 在 sources/qa_layout.md 记录。
% ================================================================

% ---------------------------------------------------------------
% 1  页面几何
%    【2026-09-28 改版】按用户指示放弃双栏，改为单栏。
%    理由：双栏每栏仅 88.5 mm，大量原理图与对照表挤在窄栏里，
%    出现多幅图元素重叠、长表超版心（实测最大 89.7 pt）。
%    单栏后版心宽度不变（183 mm），但图与表第一次真正获得整幅宽度。
%    单栏下行长变长，故把左右边距加大到 25/20 mm，使版心收窄到 165 mm
%    （10 pt 中文约 47 字/行），避免行长过长影响阅读。
%    用户同时指示「字体大小可随内容调整」，故不再强制 10 pt 正文与 8 pt 图内文字下限，
%    允许为消除重叠而按需微调；但仍不得小到影响可读性（正文 ≥9 pt、图内 ≥7 pt）。
% ---------------------------------------------------------------
\usepackage{geometry}
\geometry{
  a4paper,
  inner=25mm, outer=20mm, top=18mm, bottom=18mm,
  includeheadfoot,
  headheight=11pt, headsep=4mm, footskip=6mm
}

% ---------------------------------------------------------------
% 2  字体（企划 §6.1）
%    中文：Fandol（随 TeX Live 分发，可复现；正文宋、标题黑）
%    西文与数学：TeX Gyre Termes 及其相容数学字体（Times 风格，单一符号外观）
% ---------------------------------------------------------------
\usepackage{fontspec}
\setmainfont{TeX Gyre Termes}
\setsansfont{TeX Gyre Heros}
\setmonofont{TeX Gyre Cursor}[Scale=0.94]
\usepackage{unicode-math}
\setmathfont{TeX Gyre Termes Math}
% 【兼容性修正】unicode-math 默认把数学斜体拉丁字母映射到 U+1D400 区
% （补充平面 Mathematical Alphanumeric Symbols），例如 a → U+1D44E、t → U+1D461。
% 这类字符需要阅读器正确处理 Type0 复合字体 + 代理对；**部分中文 PDF 阅读器
% （WPS、旧版 Foxit 等）会把它渲染成实心黑块**——本项目收到过用户报告
% 「公式有黑块」，而编译期既无缺字告警、渲染器侧也复现不出。
% 实测：全书 372737 个字符里有 9605 个（2.6%）属补充平面，分布在 178/198 页上，
% 即几乎每个含公式的页面都受影响。
% 下面两行把斜体/正体拉丁字母与数字改从文本字体（TeX Gyre Termes）的
% **ASCII 码位**取字形，从而避开补充平面；希腊字母与数学符号仍用数学字体。
\setmathfont{TeX Gyre Termes Italic}[range=it/{latin,Latin}]
\setmathfont{TeX Gyre Termes}[range=up/{latin,Latin},range=sfup/{latin,Latin}]
% 希腊字母同样处理（变量里的 α、β、θ、σ 等也走补充平面 U+1D6xx 区）。
\setmathfont{TeX Gyre Termes Italic}[range=it/{greek,Greek}]
\setmathfont{TeX Gyre Termes}[range=up/{greek,Greek}]

% 10 pt 正文 / 12.5 pt 基线距。ctex 默认 15.6 pt（stretch 1.3）过松，
% 按企划 §6.1 收紧；12 pt x 1.0417 = 12.5 pt。
\linespread{1.0417}

% ---------------------------------------------------------------
% 2b 静默丢字防治（重要）
%    xeCJK 默认把 Ⅰ—ⅿ（罗马数字）与 ①—⓿（带圈字母数字）判为 Latin 类，
%    于是它们落到 TeX Gyre Termes —— 而 Termes 没有这些字形。
%    后果是**静默丢字**：不报错、不中断，只在 log 里留一行 Missing character，
%    PDF 上那个字符直接不见了。中文 Fandol 四种字体都有这些字，因此改判为 CJK 类。
%    注：U+2218 ∘ 连 Fandol 也没有，故不在此列；摄氏度改用 siunitx 的 \celsius。
% ---------------------------------------------------------------
\xeCJKDeclareCharClass{CJK}{"2160 -> "217F}   % Ⅰ—ⅿ 罗马数字
\xeCJKDeclareCharClass{CJK}{"2460 -> "24FF}   % ①—⓿ 带圈字母数字

% 数学模式里的中文：默认会被送进数学字体（TeXGyreTermesMath），该字体没有汉字，
% 于是静默丢字（实测丢过 体/床/时/积/间）。CJKmath 让数学模式中的汉字改用中文字体。
\xeCJKsetup{CJKmath=true}

% 度数符号：\circ 在 unicode-math 下映射到 U+2218，Termes 与 Fandol 都没有这个字形。
% **不要**renewcommand{\circ} —— 那会破坏 siunitx 的单位解析（实测 24 处
% "Missing { inserted"）。正确做法是给 U+2218 单独指定一个含该字形的数学字体。
\setmathfont{Latin Modern Math}[range={"2218}, Scale=MatchLowercase]

% 横向花括号：\underbrace 在 unicode-math 下用的是 U+23DF（BOTTOM CURLY BRACKET）
% 及其延长件。**TeX Gyre Termes Math 没有这组字形**，于是退化成了粗横条——
% 实测描边宽度 **4.48 pt**（正常花括号约 0.4—0.8 pt），并断成四段、中间有缺口。
% 后果：式(1-8) 的 \underbrace 处出现两条 4.5 pt 粗的黑条，用户报告为「公式有黑块」，
% 且此前一直查不到——文字层完全正常（(1)(2) 是普通 ASCII，字体已嵌入），
% **缺陷只在矢量描边里**，靠查字符/字体永远查不出来。
% 从 Latin Modern Math 补这组字形。
\setmathfont{Latin Modern Math}[range={"23DE-"23DF}, Scale=MatchLowercase]

% ---------------------------------------------------------------
% 3  数学、量纲与化学式（企划 §6.4）
% ---------------------------------------------------------------
\usepackage{amsmath}
\usepackage{mathtools}
\usepackage{siunitx}
% siunitx 2/3 兼容：两套命令名都可用，章作者不必关心版本
\providecommand{\qty}{\SI}
\providecommand{\unit}{\si}
\sisetup{
  detect-all,
  per-mode=symbol,
  range-phrase=--,
  list-units=single,
  output-product=\cdot,
  group-separator={\,},
  group-minimum-digits=5,
  % 微米符号：不要用 U+1D707（数学斜体 μ），Termes 无此字形，会静默丢字。
  % 改用直立的 U+00B5（注意 \textmu 在本环境未定义，必须用字面字符）。
  math-micro=µ,
  text-micro=µ
}
\usepackage[version=4]{mhchem}

% ---------------------------------------------------------------
% 4  图表（企划 §6.4）
% ---------------------------------------------------------------
\usepackage{graphicx}
\usepackage{xcolor}
\usepackage{needspace}
\usepackage{booktabs}
\usepackage{tabularx}
\usepackage{array}
\usepackage{multirow}
\usepackage{makecell}
\usepackage{xltabular}   % 跨页长表
\usepackage{threeparttable}
\usepackage{subcaption}
\captionsetup{
  font=small,                 % 9 pt（企划 §6.1：图表文字 8.5—9 pt）
  labelfont=bf,
  labelsep=quad,
  skip=3pt,
  justification=raggedright,
  singlelinecheck=false
}
% 题注断行兜底：题注里常有长的行内公式串，中文又缺连字点，实测会顶出版心
% （ch02-a / ch04-a / ch06-a 都出现过，最大 66 pt）。
% 注意：\captionsetup{font={small,\sloppy}} 不行 —— caption 包会校验 font 列表，
% 报 “Package caption Error: \sloppy undefined”。改为在浮动体入口处放宽。
% \@float 覆盖 figure/table，\@dblfloat 覆盖 figure*/table*。
\usepackage{etoolbox}
\makeatletter
% 【注意】\sloppy 会把 \emergencystretch 重置为 3em，比导言区全局设的 5em 还小，
% 于是「加 \sloppy 反而更容易 overfull」。所以每次 \sloppy 之后都要重设回来。
\pretocmd{\@float}{\sloppy\emergencystretch=10em}{}{\PackageWarning{dsh}{could not patch \string\@float}}
\pretocmd{\@dblfloat}{\sloppy\emergencystretch=10em}{}{\PackageWarning{dsh}{could not patch \string\@dblfloat}}
\makeatother
% 表格单元同样需要：窄列里含较长行内公式或西文串时，单元格内容会顶出列宽。
\AtBeginEnvironment{tabular}{\sloppy\emergencystretch=10em}
\AtBeginEnvironment{tabularx}{\sloppy\emergencystretch=10em}
\AtBeginEnvironment{xltabular}{\sloppy\emergencystretch=10em}
\captionsetup[sub]{font=small,labelfont=bf,labelsep=quad}

% 跨栏浮动体位置修正（twocolumn 下双栏浮动默认只能置顶）
% 注意：stfloats 与 dblfloatfix 会争抢 \c@dblbotnumber，二者只能取一。
\usepackage{dblfloatfix}
\usepackage{placeins}       % \FloatBarrier：把浮动体限制在本节内
\usepackage{flafter}        % 浮动体不得出现在其引用之前

% 图形宽度约定
% 【2026-09-28 改版】双栏 → 单栏后，图不再被挤在 88.5 mm 的窄栏里，
% 因此把 \onecolwidth 放宽到版心的约 0.8 倍（132 mm），这正是消除「图元素重叠」的主因：
% 原来标 84 mm 的图现在有 132 mm 可用，图例、轴标签、注释文字不再互相压叠。
% \twocolwidth 保留为版心满宽，供个别需要更大画幅的图使用。
% 【注意】pgfplots 的 width= 只约束坐标轴盒，轴标签、图例、长分式仍会额外外扩
% （实测过 260.2 pt > 当时栏宽 251.8 pt）。故这里仍留余量，不要直接写 \textwidth。
\newcommand{\onecolwidth}{0.80\textwidth}   % ≈132 mm
\newcommand{\twocolwidth}{0.96\textwidth}   % ≈158 mm

\usepackage{tikz}
\usepackage{pgfplots}
\pgfplotsset{compat=1.18, width=\onecolwidth}
\usetikzlibrary{arrows.meta,positioning,calc,patterns,decorations.pathreplacing,shapes.geometric,fit,backgrounds,plotmarks}
\usepgfplotslibrary{fillbetween,groupplots,colormaps}
% 图内统一字体与线宽（企划 §5.3）——灰度打印可分辨
\pgfplotsset{
  every axis/.append style={
    font=\fontsize{8.5}{10}\selectfont,
    line width=0.6pt,
    tick align=outside,
    label style={font=\fontsize{8.5}{10}\selectfont},
    tick label style={font=\fontsize{8}{9.5}\selectfont},
    legend style={font=\fontsize{8}{9.5}\selectfont, draw=none, fill=none},
    grid=none
  }
}
\tikzset{
  every node/.style={font=\fontsize{8.5}{10}\selectfont},
  thinline/.style={line width=0.5pt},
  mainline/.style={line width=0.9pt},
  axisline/.style={line width=0.7pt}
}
% 灰度安全的色彩：同时用线型/符号编码，不单靠颜色
\definecolor{accent}{gray}{0.35}
\definecolor{accentsoft}{gray}{0.68}

% 全宽（跨栏）非浮动区块：用于章首标题、宽推导、长表
% 【2026-09-28 改版】原用 cuted 的 strip 环境在双栏里做一个跨栏的章标题带。
% 单栏后整幅本来就是通栏，strip 不但无意义，在单栏下还会出错。
% 因此 fullwidth 退化为恒等环境；\chapterstart / \chapterheadstar 无需改动。
\newenvironment{fullwidth}{}{}

% ---------------------------------------------------------------
% 5  标题层级（企划 §6.1）
%    章 14 pt 黑体；一级小标题 11 pt 黑体；二级小标题 10 pt 黑体
%    最多三级；标题前 6—8 pt、后 3—4 pt；章标题不独占页
% ---------------------------------------------------------------
\usepackage{titlesec}
\titleformat{\section}[hang]
  {\heiti\fontsize{14}{17}\selectfont}
  {第\,\thesection\,章}{0.55em}{}
\titlespacing*{\section}{0pt}{7pt}{4pt}

\titleformat{\subsection}[hang]
  {\heiti\fontsize{11}{13.5}\selectfont}
  {\thesubsection}{0.5em}{}
\titlespacing*{\subsection}{0pt}{7pt}{3pt}

\titleformat{\subsubsection}[hang]
  {\heiti\fontsize{10}{12.5}\selectfont}
  {\thesubsubsection}{0.5em}{}
\titlespacing*{\subsubsection}{0pt}{6pt}{2pt}

% 第四层只作行内粗标，不编号、不进目录
\titleformat{\paragraph}[runin]
  {\heiti\normalsize}{}{0pt}{}[\,]
\titlespacing*{\paragraph}{0pt}{4pt}{0.5em}
\setcounter{secnumdepth}{3}
\setcounter{tocdepth}{2}

% 章不强制起新页（企划 §6.1「连续排版，不强制奇数页起章，不插空白页」）
% 章标题行：见下方 \chapterstar / \chapterstart / \chapterheadstar

% 不编号的章级标题（导读、综合应用、索引等）
\newcommand{\chapterheadstar}[2]{%
  \needspace{14\baselineskip}%
  \begin{fullwidth}%
    \vspace{2pt}%
    \noindent{\heiti\fontsize{14}{17}\selectfont #1}%
    \par\vspace{2.5pt}%
    \noindent{\color{accent}\rule{\linewidth}{0.8pt}}%
    \par\vspace{3pt}%
    \noindent{\fontsize{9}{11.5}\selectfont #2}%
    \par\vspace{4pt}%
  \end{fullwidth}%
  \noindent\ignorespaces
}

% ---------------------------------------------------------------
% 5b  章起始（章作者唯一入口）
%   用法：  \chapterstart{章标题}{导读与导航}
%           \label{ch:NN}
%   一次完成：编号 + 目录条目 + 页眉标记 + 全宽章标题
%   注意：章作者不要再直接使用 \section
% ---------------------------------------------------------------
\newcommand{\chapterstart}[2]{%
  \refstepcounter{section}%
  \sectionmark{#1}%
  \phantomsection
  \addcontentsline{toc}{section}{\protect\numberline{\thesection}#1}%
  \needspace{12\baselineskip}%
  \begin{fullwidth}%
    \vspace{2pt}%
    \noindent{\heiti\fontsize{14}{17}\selectfont 第\,\thesection\,章\hspace{0.55em}#1}%
    \par\vspace{2.5pt}%
    \noindent{\color{accent}\rule{\linewidth}{0.8pt}}%
    \par\vspace{3pt}%
    \noindent{\fontsize{9}{11.5}\selectfont #2}%
    \par\vspace{4pt}%
  \end{fullwidth}%
  \noindent\ignorespaces
}

% ---------------------------------------------------------------
% 6  页眉页脚（企划 §6.1）
%    页眉：章名或方法名；页脚：页码；不放装饰图案
% ---------------------------------------------------------------
\usepackage{fancyhdr}
\pagestyle{fancy}
\fancyhf{}
\renewcommand{\headrulewidth}{0.4pt}
\renewcommand{\footrulewidth}{0pt}
\fancyhead[LE]{\fontsize{8.5}{10}\selectfont 现代催化研究方法精要}
\fancyhead[RO]{\fontsize{8.5}{10}\selectfont\leftmark}
\fancyfoot[LE,RO]{\fontsize{9}{10}\selectfont\thepage}
\fancypagestyle{plain}{%
  \fancyhf{}%
  \renewcommand{\headrulewidth}{0.4pt}%
  \fancyhead[LE]{\fontsize{8.5}{10}\selectfont 现代催化研究方法精要}%
  \fancyhead[RO]{\fontsize{8.5}{10}\selectfont\leftmark}%
  \fancyfoot[LE,RO]{\fontsize{9}{10}\selectfont\thepage}%
}
% 章名写入页眉
\renewcommand{\sectionmark}[1]{\markboth{第\thesection 章\ #1}{}}

% ---------------------------------------------------------------
% 7  编号：公式、图、表按章编号（企划 §6.4）
%     形如 式(4-7)、图 4-3、表 4-2
% ---------------------------------------------------------------
\numberwithin{equation}{section}
\renewcommand{\theequation}{\thesection-\arabic{equation}}
\numberwithin{figure}{section}
\renewcommand{\thefigure}{\thesection-\arabic{figure}}
\numberwithin{table}{section}
\renewcommand{\thetable}{\thesection-\arabic{table}}

% ---------------------------------------------------------------
% 8  列表与间距（企划 §6.1：首行缩进 2 字，段间距 0—2 pt）
% ---------------------------------------------------------------
\usepackage{enumitem}
\setlist{topsep=2pt,itemsep=1pt,parsep=0pt,partopsep=0pt,leftmargin=1.8em}
\setlength{\parindent}{2em}
\setlength{\parskip}{0pt}
% 中文断行缺少可行的连字点，易产生 overfull \hbox。给 5em 应急伸缩量，
% 并放宽 tolerance：题注里常有长的行内公式串，无法断开时会顶出版心。
\setlength{\emergencystretch}{5em}
\tolerance=3000
\hbadness=10000
\setlength{\abovedisplayskip}{5pt}
\setlength{\belowdisplayskip}{5pt}
\setlength{\abovedisplayshortskip}{3pt}
\setlength{\belowdisplayshortskip}{3pt}
\setlength{\jot}{2pt}

% ---------------------------------------------------------------
% 9  链接、书签与交叉引用（企划 §6.4）
% ---------------------------------------------------------------
\usepackage{hyperref}
\hypersetup{
  unicode=true,
  bookmarksnumbered=true,
  bookmarksopen=false,
  colorlinks=false,
  pdfborder={0 0 0},
  breaklinks=true,
  pdftitle={现代催化研究方法精要：原理、读图与方法选择},
  pdfsubject={催化研究方法},
  pdfkeywords={催化；表征；谱学；动力学；方法选择}
}
\usepackage{bookmark}
\usepackage{cleveref}
\crefname{section}{章}{章}
\crefname{equation}{式}{式}
\crefname{figure}{图}{图}
\crefname{table}{表}{表}
\crefname{subsection}{节}{节}
\creflabelformat{equation}{(#2#1#3)}

% 供正文使用的章引用（避免 cleveref 在中文里的空格问题）
\newcommand{\chapref}[1]{第\,\ref{#1}\,章}
\newcommand{\equref}[1]{式\,\eqref{#1}}
\newcommand{\figref}[1]{图\,\ref{#1}}
\newcommand{\tabref}[1]{表\,\ref{#1}}
